\documentclass{jarticle}
\pagestyle{empty}
\topmargin=-2cm
\textwidth=18cm
\textheight=25.5cm
\oddsidemargin=-10mm
\evensidemargin=-10mm
\renewcommand{\baselinestretch}{1.0}
\setlength{\columnsep}{6mm}
\usepackage{color}
\newcommand{\red}[1]{\textcolor{red}{#1}}
\begin{document}

\twocolumn[
\begin{center}
\LARGE
VC発表募集（タイトル）\\
\large
\vspace{2mm}
{\hbox to 0.8\textwidth {
\mbox{画像 花子$\dagger$} \mbox{電子 太郎$\ddagger$} \mbox{学会 三郎$\ddagger$}}}
\vspace{2mm}
\begin{tabular}{ll}
$\dagger$画像大学工学部  &  $\ddagger$画像株式会社開発部
\end{tabular}\\
E-mail: $\dagger${}hanako@gazo.ac.jp, $\ddagger${}\{taro, saburo\}@denshi.co.jp
\end{center}
]

\section{はじめに}
コンピュータ・グラフィックス(CG)，及び，CGの周辺分野における革新的な技術の発表を募集します．
創造的，独創的な学術的成果をはじめ，CG技術の新しい使い方や実用例などの応用的成果まで，未来のCG業界を刺激する幅広い投稿をお待ちしております．

募集するカテゴリは，論文（査読有り，口頭発表）とポスター（査読無し）の2つです．

\red{論文に著者名，所属，謝辞などは書かないでください．WordからPDFを作成した場合，文書のプロパティに著者を特定できる情報が入ることがありますが，そのような情報は削除してください．}

\red{論文は締切前であれば何度でもアップロードできます．締切を過ぎた投稿は受け付けませんので，時間に余裕を持って投稿してください．}


\section{論文（査読有り，口頭発表）}
論文は技術論文と応用論文の2種類を募集します．

\subsection{技術論文}
学術的な研究成果について，研究の背景，理論，アルゴリズム，結果などを紹介する論文を募集します．プログラム委員会で審査を行い，採択された論文は学会当日，参加者のみに配布される予稿集に掲載されます．

\subsection{応用論文}
CG技術の新しい使い方や，ソフトウェアや映像作品での実用例の紹介など，CG技術の応用に関する論文を募集します．技術論文と同様にプログラム委員会で審査を行いますが，ここでは技術的な新規性は問わず，論文の読み易さと内容の有効性に基づいて審査を行います．採択された論文は学会当日，参加者のみに配布される予稿集に掲載されます．

\section{ポスター（査読無し）}
完成された研究のみならず，新しいアイデア，進行中の研究紹介など，さまざまな性格の発表を募集します．VCのポスターでは，「既存技術について，自分の頭で考え直したこと」もオリジナリティであるとして発表を歓迎します．優秀なポスターについてはVCポスター賞を授与します．

\small
\bibliographystyle{IEEEtran}
\nocite{*}
\bibliography{template.bib}
\end{document}
